\section{Steady-state linear response in noisy network dynamics}
% \subsection{}
We now develop a general steady-state linear response framework for noisy network dynamics under parameter perturbations. We focus on systems that relax to a stable fixed point and fluctuate around it under Gaussian white noise. When some observables change $\vec{p} \rightarrow \vec{p} + \delta \vec{p}$, the average of some observables ${\bf B}(\vec{x}; \vec{p})$ respond to this change. The goal is to express the response of steady-state observables in terms of quantities measured in the pre-perturbed steady state. We first consider the dependence of the noise-free fixed point on parameter changes. In (\ref{deterministic_force}), the noise-free fixed point satisfies
\begin{equation}
    \vec{F}(\vec{X}; \vec{p}) = 0.
\end{equation}
When the parameters change as $\vec{p} \rightarrow \vec{p} + \delta \vec{p}$, the fixed point shifts from $\vec{X}$ to $\vec{X}+\delta \vec{X}$. Expanding
\begin{equation}
    \vec{F}(\vec{X}+\delta \vec{X}; \vec{p}+ \delta \vec{p}) = 0
\end{equation}
to the linear order gives
\begin{equation}
    \nabla \vec{F}\big|_{\vec{X}; \vec{p}} \delta \vec{X} + \nabla_{\vec{p}}\vec{F}\big|_{\vec{X}; \vec{p}} \delta \vec {p} = 0.
\end{equation}
With $\mathbf{Q}\equiv\nabla\vec{F}\big|_{\vec{X}}$, and assuming $\mathbf{Q}^{-1}$ exists, one obtains
\begin{equation}
    \frac{\partial X_i}{\partial p_\alpha} = -\left(\mathbf{Q}^{-1} \nabla_{\vec{p}} \vec{F}\big|_{\vec{X}; \vec{p}}\right)_{i\alpha}.
\end{equation}
This relation shows that the shift of the fixed point induced by a parameter perturbation is controlled by the linearized dynamics through $\mathbf{Q}$. For an observable $\mathbf{B}(\vec{x}; \vec{p})$, the steady-state linear response is defined by
\begin{equation}
    \langle \mathbf{B} \rangle_{ss, \vec{p}+\delta \vec{p}} = \langle \mathbf{B} \rangle_{ss, \vec{p}} + \boldsymbol{\chi} \delta \vec{p}.
    \label{response}
\end{equation}
The steady-state response function is then
\begin{equation}
    \boldsymbol{\chi} = \langle \nabla_{\vec{p}} \mathbf{B} \rangle_{ss, \vec{p}} + \int d\vec{x} \mathbf{B} \nabla_{\vec{p}} \psi_{ss}.
    \label{response_function}
\end{equation}
Defining the conjugate force with respect to parameters
\begin{equation}
    \vec{\mathscr{F}} \equiv -\nabla_{\vec{p}} \Phi,
    \label{conjugate_force}
\end{equation} 
one obtains
\begin{equation}
    \boldsymbol{\chi} = \langle \nabla_{\vec{p}} \mathbf{B} \rangle_{ss, \vec{p}} + \langle \delta \mathbf{B} \delta \vec{\mathscr{F}}^\intercal \rangle_{ss, \vec{p}}.
    \label{response_function2}
\end{equation}
This is the general steady-state linear response formula used throughout this work. Importantly, it does not require equilibrium. Instead, it only requires that the system fluctuates around a steady state. In the following sections, we apply this general framework to parameter changes in the intrinsic node parameter $r_\alpha$ and the connection weight $W_{\alpha\beta}$.
